\section*{Hoofdstuk \ref{ch:b3:total-synthesis} --- Totaalsynthese: strategie en klassieke routes}

\begin{solution}{exo:b3:total-synthesis:1}
Bij 90\,\%: $0.9^{10} = 35$\,\% en $0.9^{20} = 12$\,\%. Bij 80\,\%: $0.8^{10} = 11$\,\% en $0.8^{20} = 1.2$\,\%.
\end{solution}

\begin{solution}{exo:b3:total-synthesis:2}
$5 + 3 + 1 + 4 = 13$ stappen in totaal; LLS $5 + 1 + 4 = 10$.
\end{solution}

\begin{solution}{exo:b3:total-synthesis:3}
4-(Dimethylamino)butan-2-on, $\ce{CH3COCH2CH2N(CH3)2}$: het enol van aceton addeert aan het
iminiumion $\ce{CH2=N+(CH3)2}$.
\end{solution}

\begin{solution}{exo:b3:total-synthesis:4}
Disconnecteer de twee ringbindingen die allylisch staan ten opzichte van
de C=C: buta-1,3-dieen en but-3-een-2-on
(methylvinylketon), die bij verhitting reageren.
\end{solution}

\begin{solution}{exo:b3:total-synthesis:5}
Lineair: $0.9^{21} = 11$\,\%. Convergent: $0.9^{11} = 31$\,\%, bijna drie keer zoveel.
\end{solution}

\begin{solution}{exo:b3:total-synthesis:6}
A: $(7 + 1)/12 = 67$\,\%. B: $7/9 = 78$\,\%.
\end{solution}

\begin{solution}{exo:b3:total-synthesis:7}
Zonder: $0.9^8 = 43$\,\%. Met vier extra stappen bij 95\,\%: $43\,\% \times 0.95^4 = 35$\,\%.
\end{solution}

\begin{solution}{exo:b3:total-synthesis:8}
Primair: een omvangrijke silylether (\textit{tert}-butyldifenylsilyl), verwijderd met fluoride.
Secundair: een 4-methoxybenzylether, verwijderd door oxidatie met DDQ. De
tertiaire alcohol, gehinderd, kan vaak vrij blijven, of beschermd worden
als een kleine silylether die door mild zuur gesplitst wordt. Elke
omstandigheid laat de andere groepen ongemoeid.
\end{solution}

\begin{solution}{exo:b3:total-synthesis:9}
De \omterm{def:b3:total-synthesis:total}{totaalsyntheses} vereisten tientallen stappen met \omterm{def:b3:total-synthesis:architecture}{totaalrendementen} ver
onder 1\,\%; een verbinding met het volledige ringsysteem en de meeste
stereocentra, gewonnen uit de naalden van de Europese taxus, een
hernieuwbare bron, wordt in enkele stappen omgezet.
\end{solution}

\begin{solution}{exo:b3:total-synthesis:10}
Michaeladditie: het enolaat van het dion (C2, tussen de carbonylgroepen) addeert aan het
$\beta$-koolstofatoom van methylvinylketon, wat een triketon geeft. Aldol: het enolaat
van het methylketon valt intramoleculair een carbonylgroep van de ring
aan, en dehydratatie geeft het cyclohexenon: het keton van
Wieland-Miescher, een bicyclisch enedion met een angulaire methylgroep.
\end{solution}

\begin{solution}{exo:b3:total-synthesis:11}
$y^{m+1}/y^{2m+1} = y^{-m}$: bij 90\,\% 1.7 voor $m = 5$ en 4.9 voor $m = 15$. Hoe langer de route, hoe meer
convergentie loont.
\end{solution}

\begin{solution}{exo:b3:total-synthesis:12}
Elke ring sluit met het kation en de aanvallende dubbele binding in een
stoelvormige schikking; de additie over elke dubbele binding is anti, en
de substituenten nemen pseudo-equatoriale posities in, wat elke
ringverbinding \textit{trans} maakt. Een dubbele binding \textit{E}
plaatst haar twee ketenvoortzettingen zo dat de \textit{trans}-verbinding ontstaat; een binding \textit{Z} zou
de \textit{cis}-verbinding geven: de geometrie van het polyeen codeert de stereochemie van
het product.
\end{solution}

\begin{solution}{pb:b3:total-synthesis:1}
\textbf{1.} $\ce{C4H6O2}$, $\ce{CH5N}$, $\ce{C5H6O5}$, $\ce{C8H13NO}$.

\textbf{2.} $\ce{C4H6O2 + CH5N + C5H6O5 -> C8H13NO + 2CO2 + 2H2O}$.

\textbf{3.} 86.0, 31.0 en \qty{146.0}{g/mol}; tropinon \qty{139.0}{g/mol}.

\textbf{4.} $139.0/(86.0 + 31.0 + 146.0) = 0.53$: 53\,\%.

\textbf{5.} Zijn centrale groepen $\ce{CH2}$, elk tussen twee carbonylgroepen, enoliseren gemakkelijk in
water; aceton zelf is te zwak nucleofiel, en de activerende
carboxylgroepen vertrekken achteraf.

\textbf{6.} De iminiumionen hebben vrij amine nodig (te veel zuur
protoneert het), en het enol en de vorming van het iminiumion vereisen
wat zuur: een pH rond neutraal brengt ze in evenwicht.

\textbf{7.} $\ce{CH3NH2}$ addeert aan één aldehydegroep; na verlies van water
ontstaat het iminiumion
$\ce{R-CH=N+H-CH3}$.

\textbf{8.} Een enolkoolstofatoom van acetondicarbonzuur valt het
iminiumkoolstofatoom aan, wat een binding C--C en een secundair amine
vormt.

\textbf{9.} Het secundaire amine condenseert met de tweede aldehydegroep
van dezelfde keten, wat een cyclisch iminiumion (vijfring) geeft.

\textbf{10.} Het enol aan de andere kant van het keton valt het aan, wat
de zesring en het bicyclische skelet sluit.

\textbf{11.} Elk is een $\beta$-ketozuur, dat $\ce{CO2}$ verliest via een cyclische zesledige
overgangstoestand, wat het enol van het keton geeft.

\textbf{12.} Twee bindingen C--C en twee bindingen C--N.

\textbf{13.} $0.75^{14} = 0.018$: 1.8\,\%.

\textbf{14.} $42/1.78 = 24$.

\textbf{15.} 100\,\%: één stap, allemaal opbouw van bindingen.

\textbf{16.} Eén bewerking, water als oplosmiddel, goedkope componenten,
geen beschermende groepen en bijna geen afval behalve koolstofdioxide en
water.

\textbf{17.} De plant bouwt het tropaanskelet op uit een van een aminozuur
afgeleid cyclisch iminiumion en een van acetaat afgeleide eenheid, via
hetzelfde soort mannichchemie.

\textbf{18.} Cascades (en ringvormende \omterm{def:b3:pericyclic:families}{cycloaddities} zoals de
diels-alderreactie).

\textbf{19.} Nee: een \omterm{def:b3:point-groups:mirror}{spiegelvlak} gaat door N, zijn methylgroep, C3 en het
carbonylzuurstofatoom.

\textbf{20.} Elk heeft vier verschillende groepen, maar ze zijn
spiegelbeelden van elkaar in het eigen symmetrievlak van het molecuul:
een verbinding van het mesotype.

\textbf{21.} Diastereomeren, die verschillen in de oriëntatie van de OH op
C3 ten opzichte van de stikstofbrug.

\textbf{22.} Nee: het \omterm{def:b3:point-groups:mirror}{spiegelvlak} blijft in beide behouden.

\textbf{23.} De twee vlakken van de carbonylgroep zijn niet gelijkwaardig:
het ene is naar de stikstofbrug gericht, het andere naar de brug van
twee koolstofatomen, en die hinderen verschillend.

\textbf{24.} De eenpotsroute geeft ongeveer 24 keer het \omterm{def:b3:total-synthesis:architecture}{totaalrendement} van de lineaire route
van 14 stappen.
\end{solution}
